% =============================================================================
% APPENDIX J: LIT-RECENT: Ten Recent Papers (2024-2025)
% =============================================================================
% Module: EBF_J_LIT_RECENT
% Version: 2.0 (January 2026)
% Status: Draft
% Language: English
%
% This appendix demonstrates the EBF framework's prospective application
% by analyzing ten papers published in 2024-2025. Unlike retrospective
% Nobel analysis, these papers were published after the framework was
% conceived, testing whether any new economics paper can be characterized
% in terms of $(C, \Psi, C^*, K, Q)$.
%
% Category: LIT
% Core Question: Can the EBF framework successfully characterize new
%                economics research published after the framework's development?
%
% =============================================================================

% =============================================================================
% CROSS-REFERENCE STRUCTURE
% =============================================================================
%
% CHAPTER LINKAGE (Required):
% - Primary Chapter: Chapter 13: Literature Integration
% - Secondary Chapters: Chapter 2 (Framework Overview), Chapter 15 (Future)
%
% This appendix provides prospective validation of framework generality:
%
% DEPENDENCIES (what this appendix REQUIRES):
% - Appendix HOW (CORE-HOW): Complementarity framework ($\gamma$)
% - Appendix CTW (CORE-WHEN): Context framework ($\Psi$)
% - Appendix WAT (CORE-WHAT): Welfare dimensions (FEPSDE)
% - Appendix NOB (LIT-NOBEL): Retrospective Nobel integration
%
% DEPENDENTS (what REQUIRES this appendix):
% - Chapter 15: Uses prospective analysis for future directions
% - Other LIT appendices: Reference methodology for paper integration
%
% =============================================================================

\begin{tcolorbox}[colback=purple!5!white, colframe=purple!75!black,
    title=Chapter Linkage]

\textbf{Primary Chapter:} Chapter 13: Literature Integration
\begin{itemize}[nosep]
\item Section 13.1: Integration Methodology $\leftrightarrow$ This Appendix Overview
\item Section 13.4: Prospective Application $\leftrightarrow$ This Appendix Analysis
\end{itemize}

\textbf{Secondary Chapters:}
\begin{itemize}[nosep]
\item Chapter 2: Framework Overview --- validates framework generality
\item Chapter 15: Future Directions --- identifies research frontiers
\end{itemize}

\textbf{Reading Guide:}
\begin{itemize}[nosep]
\item For conceptual overview: Read Chapter 13 first
\item For retrospective analysis: See Appendix NOB (LIT-NOBEL)
\item For specific domains: See relevant DOMAIN appendices
\end{itemize}

\end{tcolorbox}

%%%%%%%%%%%%%%%%%%%%%%%%%%%%%%%%%%%%%%%%%%%%%%%%%%%%%%%%%%%%%%%%%%%%%%%%%%%%%%%
\chapter{LIT-RECENT: Ten Recent Papers (2024--2025)}
\label{app:recentpapers}

\begin{tcolorbox}[colback=blue!5!white, colframe=blue!75!black,
    title=Quick Reference: Appendix REC]

\textbf{Appendix:} REC (LIT)

\textbf{Core Question:} What are the key findings in this domain?

\textbf{Key Concepts:}
\begin{itemize}[nosep]
\item Framework integration
\item Parameter validation
\item Cross-references
\end{itemize}

\textbf{Cross-References:} See related appendices below
\end{tcolorbox}

%%%%%%%%%%%%%%%%%%%%%%%%%%%%%%%%%%%%%%%%%%%%%%%%%%%%%%%%%%%%%%%%%%%%%%%%%%%%%%%

\begin{abstract}
This appendix demonstrates the EBF framework's \emph{prospective} application
by analyzing ten papers published in 2024--2025. Unlike the Nobel retrospective
(Appendix NOB), these papers were published \emph{after} the framework was
conceived. This tests the claim that any new economics paper can be
characterized in terms of $(C, \Psi, C^*, K, Q)$. All ten papers map
successfully, with complementarities ($C^{SS}$, $C^{market}$, $C^{F,Political}$),
context dimensions ($\Psi_{neighborhood}$, $\Psi_{policy}$, $\Psi_{inst}$),
and all FEPSDE welfare dimensions appearing across the sample.
\end{abstract}

\begin{tcolorbox}[colback=blue!5!white, colframe=blue!75!black,
    title=Quick Reference: Recent Papers Integration]

\textbf{Core Question:} Can the EBF framework successfully characterize new
economics research published after the framework's development?

\textbf{Key Result:}
\begin{equation}
\text{All 10 papers} \mapsto (C, \Psi, C^*, K, Q) \quad \checkmark
\end{equation}

\textbf{Key Concepts:}
\begin{itemize}[nosep]
\item Prospective validation: Papers published 2024-2025, after framework development
\item Framework elements: Complementarities, context, equilibria, coherence, welfare
\item Restrictions revealed: Framework shows what each paper does \emph{not} model
\item Cross-domain coverage: Social capital, climate, banking, trade, gender, AI
\end{itemize}

\textbf{Cross-References:} Appendix NOB (LIT-NOBEL), Appendix HOW (CORE-HOW),
Appendix CTW (CORE-WHEN), Appendix WAT (CORE-WHAT)
\end{tcolorbox}


%%%%%%%%%%%%%%%%%%%%%%%%%%%%%%%%%%%%%%%%%%%%%%%%%%%%%%%%%%%%%%%%%%%%%%%%%%%%%%%
% -----------------------------------------------------------------------------
% APPENDIX SCOPE BOX
% -----------------------------------------------------------------------------

\begin{tcolorbox}[
    colback=orange!5!white,
    colframe=orange!75!black,
    title={\textbf{Appendix Scope: Ziel / In-Scope / Out-of-Scope / Constraints}},
    fonttitle=\bfseries
]

\textbf{Ziel (Objective):}
\begin{itemize}[nosep]
    \item Can the EBF framework $(C, \Psi, C^*, K, Q)$ successfully characterize
economics papers published in 2024--2025, validating its prospective
integration capability?
\end{itemize}

\textbf{In-Scope:}
\begin{itemize}[nosep]
    \item The Fundamental Question
    \item Overview: Ten Recent Papers Mapped
    \item Detailed Integration
    \item Key Results: Synthesis
    \item Worked Example
\end{itemize}

\textbf{Out-of-Scope (delegiert an andere Appendices):}
\begin{itemize}[nosep]
    \item CORE-HOW $\rightarrow$ Appendix HOW
    \item CORE-WHEN $\rightarrow$ Appendix CTW
    \item CORE-WHAT $\rightarrow$ Appendix WAT
    \item LIT-NOBEL $\rightarrow$ Appendix NOB
    \item LIT-NOBEL $\rightarrow$ Appendix NOB
\end{itemize}

\textbf{Constraints (Anwendungsgrenzen):}
\begin{itemize}[nosep]
    \item [TODO: Define when this appendix does NOT apply]
    \item [TODO: Define required prerequisites]
    \item [TODO: Define known limitations]
\end{itemize}

\textbf{Lieferobjekte (Deliverables):}
\begin{itemize}[nosep]
    \item 1 Table(s)
    \item 3 Equation(s)
    \item Worked Example(s)
\end{itemize}

\end{tcolorbox}


\section{The Fundamental Question}
\label{sec:J-fundamental}
%%%%%%%%%%%%%%%%%%%%%%%%%%%%%%%%%%%%%%%%%%%%%%%%%%%%%%%%%%%%%%%%%%%%%%%%%%%%%%%

\subsection{Why This Matters}

A framework's value depends on its ability to integrate not just past research
but future research as well. While retrospective integration (Appendix NOB)
demonstrates that EBF can unify established findings, prospective application
is the stronger test:

\begin{itemize}
    \item \textbf{Retrospective:} Can EBF explain what we already knew?
    \item \textbf{Prospective:} Can EBF characterize what we are learning now?
\end{itemize}

If EBF can successfully integrate papers published after its development,
this provides strong evidence for its generality and continued relevance.

\subsection{The Core Problem}

The challenge is demonstrating that:
\begin{enumerate}
    \item Papers from diverse subfields all map to $(C, \Psi, C^*, K, Q)$
    \item The mapping reveals substantive insights, not just relabeling
    \item Framework identifies what each paper assumes vs. endogenizes
    \item Integration suggests productive research extensions
\end{enumerate}

\begin{tcolorbox}[colback=red!5!white, colframe=red!75!black,
    title=The Fundamental Question]
\textbf{Can the EBF framework $(C, \Psi, C^*, K, Q)$ successfully characterize
economics papers published in 2024--2025, validating its prospective
integration capability?}
\end{tcolorbox}


%%%%%%%%%%%%%%%%%%%%%%%%%%%%%%%%%%%%%%%%%%%%%%%%%%%%%%%%%%%%%%%%%%%%%%%%%%%%%%%
\section{Overview: Ten Recent Papers Mapped}
\label{sec:J-overview}
%%%%%%%%%%%%%%%%%%%%%%%%%%%%%%%%%%%%%%%%%%%%%%%%%%%%%%%%%%%%%%%%%%%%%%%%%%%%%%%

\begin{center}
\renewcommand{\arraystretch}{1.15}
\small
\begin{tabular}{lll}
\hline
\textbf{Paper} & \textbf{Topic} & \textbf{Framework Elements} \\
\hline
Chetty et al. (2024) & Social capital, mobility & $C^{SS}$, $\Psi_{neighborhood}$ \\
Stantcheva et al. (2025) & Climate policy attitudes & $U^{Eco}$, $U^F$, $\Psi_{info}$ \\
Amador \& Bianchi (2024) & Bank runs, credit easing & $K_{financial}$, $\Psi_{policy}$ \\
Choi et al. (2024) & NAFTA local effects & $C^{trade}$, $\Psi_{inst}^{trade}$ \\
Fenizia \& Saggio (2024) & Mafia and growth & $\Psi_{inst}^{extractive}$, $Q$ \\
Kwon et al. (2024) & Corporate concentration & $C_{ij}^{market}$, $\Psi_{tech}$ \\
Herbst \& Hendren (2024) & Human capital financing & $C^{FP}$, info asymmetry \\
AI \& Inequality (2024) & AI and wealth gap & $\Psi_{tech}^{AI}$, $U^F$ distribution \\
Goldin (2024) & Gender economics (Nobel) & $C^*$ by gender, $\Psi_{norm}$ \\
Digital Distraction (2025) & Mobile apps, outcomes & $\Psi_{digital}$, $U^D$, $U^F$ \\
\hline
\end{tabular}
\end{center}


%%%%%%%%%%%%%%%%%%%%%%%%%%%%%%%%%%%%%%%%%%%%%%%%%%%%%%%%%%%%%%%%%%%%%%%%%%%%%%%
\section{Detailed Integration}
\label{sec:J-detailed}
%%%%%%%%%%%%%%%%%%%%%%%%%%%%%%%%%%%%%%%%%%%%%%%%%%%%%%%%%%%%%%%%%%%%%%%%%%%%%%%

%--- Paper 1 ---
\subsection{Paper 1: Chetty et al. (2024) -- Social Capital and Economic Mobility}

\paragraph{Citation.}
Chetty, R., Jackson, M.O., Kuchler, T., Stroebel, J. et al. (2022/2024).
``Social Capital I \& II.'' \emph{Nature} and \emph{QJE}.

\paragraph{Core Finding.}
Cross-class social connections (``economic connectedness'') strongly predict
upward mobility.

\paragraph{Framework Integration.}
\begin{itemize}
  \item \textbf{Complementarity:} $C^{SS}_{low,high} > 0$---cross-class connections
        complement mobility
  \item \textbf{Context:} $\Psi_{neighborhood}$ determines connection opportunities
  \item \textbf{Welfare:} Affects $U^F$ (income) and $U^S$ (social integration)
  \item \textbf{Coherence:} Segregated communities: $K^{cross-class} \approx 0$
\end{itemize}

\paragraph{Restrictions Assumed.}
$\Psi_{neighborhood}$ largely exogenous; $U^{IDN}$ effects underdeveloped.

\begin{tcolorbox}[colback=yellow!5!white, colframe=yellow!75!black,
    title=EBF Insight: Chetty et al.]
The paper identifies a specific complementarity ($C^{SS}_{low,high}$) and its
context dependence ($\Psi_{neighborhood}$). EBF suggests extensions: How does
$\Psi_{norm}$ (social norms) affect willingness to form cross-class ties?
\end{tcolorbox}


%--- Paper 2 ---
\subsection{Paper 2: Stantcheva et al. (2025) -- Climate Policy Attitudes}

\paragraph{Citation.}
Dechezleprêtre et al. (2025). ``Fighting Climate Change: International Attitudes.''
\emph{AER}, 115(4).

\paragraph{Core Finding.}
Support for climate policies depends on perceived effectiveness and fairness,
\emph{not} climate concern itself.

\paragraph{Framework Integration.}
\begin{itemize}
  \item \textbf{Welfare trade-off:} $U^{Eco}$ vs. $U^F$
  \item \textbf{Context:} $\Psi_{information}$ about policy effectiveness
  \item \textbf{Cross-domain:} $C^{Eco,F}$ perceived as negative (short-run substitutes)
\end{itemize}

\paragraph{Framework Insight.}
``Concern doesn't predict support'' = acceptance depends on $C^{Eco,F}$ perception,
not $U^{Eco}$ level.

\begin{tcolorbox}[colback=yellow!5!white, colframe=yellow!75!black,
    title=EBF Insight: Climate Policy]
The finding that concern $\neq$ support reveals that perceived complementarity
between environmental and financial welfare ($C^{Eco,F}$) mediates policy
acceptance. EBF predicts that reframing to emphasize co-benefits would increase
support.
\end{tcolorbox}


%--- Paper 3 ---
\subsection{Paper 3: Amador \& Bianchi (2024) -- Bank Runs and Credit Easing}

\paragraph{Citation.}
\emph{AER}, 114(7): 2073-2110.

\paragraph{Core Finding.}
Central bank policy prevents runs by changing depositor expectations.

\paragraph{Framework Integration.}
This is a textbook $(C, \Psi, C^*)$ paper:
\begin{equation}
  C^*_{no\ run}(\Psi_{CB\ credible}) \neq C^*_{run}(\Psi_{CB\ not\ credible})
\end{equation}
\begin{itemize}
  \item \textbf{Multiple equilibria:} ``Run'' and ``No run'' both stable
  \item \textbf{Policy as context:} $\Psi_{policy}^{CB}$ selects equilibrium
  \item \textbf{Coherence collapse:} Bank run is $K_{bank} \to 0$
\end{itemize}

\begin{tcolorbox}[colback=yellow!5!white, colframe=yellow!75!black,
    title=EBF Insight: Bank Runs]
Perfect illustration of how policy context ($\Psi_{policy}$) determines
which equilibrium ($C^*$) is selected. Links directly to Appendix CIA
(DOMAIN-CAPITAL) on coordination failures.
\end{tcolorbox}


%--- Paper 4 ---
\subsection{Paper 4: Choi et al. (2024) -- Local Effects of NAFTA}

\paragraph{Citation.}
\emph{AER}, 114(6): 1540-1575.

\paragraph{Core Finding.}
NAFTA's economic shocks had lasting political effects.

\paragraph{Framework Integration.}
\begin{itemize}
  \item \textbf{Cross-domain spillover:} $C^{F \to Political} > 0$
  \item \textbf{Context hysteresis:} $\Psi_{political}$ shifted permanently
  \item \textbf{Identity:} $U^{IDN}_{manufacturing}$ affected by job loss
\end{itemize}

\paragraph{Prediction Validated.}
Evidence for framework's ``cross-domain spillover with lags'' prediction.

\begin{tcolorbox}[colback=yellow!5!white, colframe=yellow!75!black,
    title=EBF Insight: Trade and Politics]
Economic shocks ($\Delta U^F$) spill over to political preferences
($C^{F \to Political}$), validating EBF's cross-domain complementarity
structure. The persistence ($\Psi$ hysteresis) shows context is sticky.
\end{tcolorbox}


%--- Paper 5 ---
\subsection{Paper 5: Fenizia \& Saggio (2024) -- Mafia and Growth}

\paragraph{Citation.}
\emph{AER}, 114(7): 2171-2200.

\paragraph{Core Finding.}
Mafia infiltration reduces growth by 15-20\%.

\paragraph{Framework Integration.}
\begin{itemize}
  \item \textbf{Extractive institutions:} $\Psi_{inst}^{mafia}$
  \item \textbf{Coherence trap:} High $K$ (stable) but low $Q$ (bad)
  \item \textbf{$\Psi_{inst} \to Q$:} Direct Acemoglu test
\end{itemize}

\begin{tcolorbox}[colback=yellow!5!white, colframe=yellow!75!black,
    title=EBF Insight: Extractive Institutions]
Illustrates how high coherence ($K$) does not guarantee good outcomes ($Q$).
Mafia-controlled regions are stable but inefficient---a ``bad equilibrium''
that demonstrates the distinction between $K$ and $Q$.
\end{tcolorbox}


%--- Paper 6 ---
\subsection{Paper 6: Kwon et al. (2024) -- Corporate Concentration}

\paragraph{Citation.}
\emph{AER}, 114(7): 2111-2140.

\paragraph{Core Finding.}
Concentration rose steadily since 1920, driven by technology.

\paragraph{Framework Integration.}
\begin{itemize}
  \item \textbf{Technology context:} $\Psi_{tech}$ drives concentration
  \item \textbf{Market complementarity:} Increasing returns $\Rightarrow C_{ij}^{market} > 0$
  \item \textbf{Long-run dynamics:} $\Psi_{tech}$ evolved over 100 years
\end{itemize}

\begin{tcolorbox}[colback=yellow!5!white, colframe=yellow!75!black,
    title=EBF Insight: Market Structure]
Technology context ($\Psi_{tech}$) shapes market complementarity structure,
which in turn determines concentration. Long time horizon shows slow-moving
context effects.
\end{tcolorbox}


%--- Paper 7 ---
\subsection{Paper 7: Herbst \& Hendren (2024) -- Human Capital Financing}

\paragraph{Citation.}
\emph{AER}, 114(7): 2024-2072.

\paragraph{Core Finding.}
Information asymmetries prevent efficient human capital investment.

\paragraph{Framework Integration.}
\begin{itemize}
  \item \textbf{Cross-domain:} $C^{FP} > 0$---financial and human capital complements
  \item \textbf{Information in context:} $\Psi_{information}$ determines achievable $C^*$
  \item \textbf{Market failure:} Efficient $C^*$ unreachable
\end{itemize}

\begin{tcolorbox}[colback=yellow!5!white, colframe=yellow!75!black,
    title=EBF Insight: Human Capital]
Information asymmetry ($\Psi_K$) prevents reaching efficient equilibrium.
The complementarity between financial and personal capital ($C^{FP}$)
means underinvestment in one reduces returns to the other.
\end{tcolorbox}


%--- Paper 8 ---
\subsection{Paper 8: AI and Wealth Inequality (2024)}

\paragraph{Citation.}
Multiple 2024 papers on AI and inequality.

\paragraph{Core Finding.}
AI adoption correlates with increased wealth inequality.

\paragraph{Framework Integration.}
\begin{itemize}
  \item \textbf{Skill complementarity:} $C^{AI,high-skill} > 0$, $C^{AI,routine} < 0$
  \item \textbf{Context speed mismatch:} $\Psi_{tech}^{AI}$ moving faster than $\Psi_{inst}$
\end{itemize}

\paragraph{Framework Prediction.}
Coherence problems when context components move at different speeds.

\begin{tcolorbox}[colback=yellow!5!white, colframe=yellow!75!black,
    title=EBF Insight: AI and Inequality]
AI creates asymmetric complementarities---positive for high-skill workers,
negative for routine workers. The mismatch between fast-moving $\Psi_{tech}$
and slow-moving $\Psi_{inst}$ creates transition frictions.
\end{tcolorbox}


%--- Paper 9 ---
\subsection{Paper 9: Goldin (2024) -- Nobel Lecture on Gender}

\paragraph{Citation.}
\emph{AER}, 114(6): 1515-1539.

\paragraph{Core Finding.}
Persistent gender gap explained by ``greedy jobs.''

\paragraph{Framework Integration.}
\begin{itemize}
  \item \textbf{Gendered $C^*$:} Historical $C^*_{male} \neq C^*_{female}$
  \item \textbf{Structural barrier:} $C^{work,family} < 0$ for ``greedy jobs''
  \item \textbf{Coherence problem:} Labor market $C^*$ conflicts with family $C^*$
\end{itemize}

\begin{tcolorbox}[colback=yellow!5!white, colframe=yellow!75!black,
    title=EBF Insight: Gender Economics]
``Greedy jobs'' represent negative complementarity between work and family
($C^{work,family} < 0$). The gender gap persists because this structural
feature interacts with gendered social norms ($\Psi_{norm}$).
\end{tcolorbox}


%--- Paper 10 ---
\subsection{Paper 10: Digital Distraction (2025)}

\paragraph{Citation.}
QJE paper on mobile apps and outcomes.

\paragraph{Core Finding.}
Mobile usage affects academic/labor outcomes with peer effects.

\paragraph{Framework Integration.}
\begin{itemize}
  \item \textbf{Digital welfare:} $U^D$ as distinct dimension
  \item \textbf{Cross-domain:} $C^{DF} < 0$---digital may substitute for $U^F$
  \item \textbf{Peer complementarity:} $C^{D}_{peer} > 0$---usage is contagious
\end{itemize}

\begin{tcolorbox}[colback=yellow!5!white, colframe=yellow!75!black,
    title=EBF Insight: Digital Behavior]
Digital dimension ($U^D$) interacts with other welfare dimensions through
complementarities. Peer effects ($C^D_{peer} > 0$) amplify individual choices.
Links to FEPSDE dimension structure in Appendix WAT.
\end{tcolorbox}


%%%%%%%%%%%%%%%%%%%%%%%%%%%%%%%%%%%%%%%%%%%%%%%%%%%%%%%%%%%%%%%%%%%%%%%%%%%%%%%
\section{Key Results: Synthesis}
\label{sec:J-results}
%%%%%%%%%%%%%%%%%%%%%%%%%%%%%%%%%%%%%%%%%%%%%%%%%%%%%%%%%%%%%%%%%%%%%%%%%%%%%%%

\subsection{All Ten Papers Map Successfully}

Every paper characterized using $(C, \Psi, C^*, K, Q)$:

\begin{center}
\renewcommand{\arraystretch}{1.3}
\begin{tabular}{|l|l|l|l|}
\hline
\textbf{Paper} & \textbf{Complementarity} & \textbf{Context} & \textbf{Welfare} \\
\hline
Chetty et al. & $C^{SS}_{low,high}$ & $\Psi_{neighborhood}$ & $U^F$, $U^S$ \\
Stantcheva et al. & $C^{Eco,F}$ & $\Psi_{information}$ & $U^{Eco}$, $U^F$ \\
Amador \& Bianchi & $C^*_{run/no-run}$ & $\Psi_{policy}^{CB}$ & $U^F$ \\
Choi et al. & $C^{F \to Political}$ & $\Psi_{political}$ & $U^F$, $U^{IDN}$ \\
Fenizia \& Saggio & $C^{inst}$ & $\Psi_{inst}^{mafia}$ & $Q$ aggregate \\
Kwon et al. & $C_{ij}^{market}$ & $\Psi_{tech}$ & $U^F$ \\
Herbst \& Hendren & $C^{FP}$ & $\Psi_K$ (asymmetry) & $U^F$, $U^P$ \\
AI \& Inequality & $C^{AI,skill}$ & $\Psi_{tech}^{AI}$ & $U^F$ distribution \\
Goldin & $C^{work,family}$ & $\Psi_{norm}$ & $U^F$, $U^S$ \\
Digital Distraction & $C^{DF}$, $C^D_{peer}$ & $\Psi_{digital}$ & $U^D$, $U^F$ \\
\hline
\end{tabular}
\end{center}

\subsection{Restrictions Made Explicit}

The framework reveals what each paper does \emph{not} model---not as criticism
but as identification of extension opportunities:

\begin{itemize}
    \item \textbf{Chetty et al.:} $\Psi_{neighborhood}$ exogenous; could endogenize
    \item \textbf{Stantcheva et al.:} $\Psi_{norm}$ not analyzed; could add
    \item \textbf{Amador \& Bianchi:} $U^{IDN}$ (identity) not modeled
    \item \textbf{Choi et al.:} $K$ dynamics not formalized
    \item \textbf{Fenizia \& Saggio:} $C^*$ transition paths not specified
\end{itemize}

\subsection{Prospective Claim Supported}

These papers---published 2024-2025---can all be characterized in framework terms.
The framework can integrate new research as it appears.

\begin{tcolorbox}[colback=green!5!white, colframe=green!75!black,
    title=Prospective Validation Result]
\textbf{All 10 papers map to $(C, \Psi, C^*, K, Q)$}

This demonstrates that EBF is not merely an ex-post rationalization of known
findings, but a genuine organizing framework capable of integrating ongoing
research across behavioral economics, finance, trade, gender, digital, and
institutional economics.
\end{tcolorbox}


%%%%%%%%%%%%%%%%%%%%%%%%%%%%%%%%%%%%%%%%%%%%%%%%%%%%%%%%%%%%%%%%%%%%%%%%%%%%%%%
\section{Worked Example}
\label{sec:J-example}
%%%%%%%%%%%%%%%%%%%%%%%%%%%%%%%%%%%%%%%%%%%%%%%%%%%%%%%%%%%%%%%%%%%%%%%%%%%%%%%

\subsection{Integrating a New Paper: Step-by-Step Method}

\paragraph{Setup.}
We demonstrate the method for integrating new papers into EBF using a
hypothetical 2026 paper on ``Remote Work and Urban Housing Markets.''

\paragraph{Application: Five-Step Process.}

\begin{enumerate}
\item \textbf{Identify the core finding:}
    ``Remote work adoption reduces urban housing demand, with heterogeneous
    effects by occupation and amenity value.''

\item \textbf{Map to complementarity structure:}
    \begin{itemize}
    \item $C^{remote, urban} < 0$: Remote work substitutes for urban location
    \item $C^{amenity, urban} > 0$: Amenities complement urban location
    \item Net effect depends on relative magnitudes
    \end{itemize}

\item \textbf{Identify context dimensions:}
    \begin{itemize}
    \item $\Psi_{tech}$: Technology enabling remote work
    \item $\Psi_{norm}$: Employer acceptance of remote work
    \item $\Psi_{occupation}$: Occupation-specific remote feasibility
    \end{itemize}

\item \textbf{Characterize equilibrium:}
    \begin{equation}
    C^*_{housing}(\Psi_{remote}) = f(C^{remote,urban}, C^{amenity,urban}, \Psi)
    \end{equation}
    Multiple equilibria possible: high-density urban vs. distributed.

\item \textbf{Note restrictions and extensions:}
    \begin{itemize}
    \item Paper assumes $\Psi_{tech}$ exogenous
    \item Does not model $K$ dynamics (how fast cities adjust)
    \item Extension: Add $U^S$ (social welfare of urban density)
    \end{itemize}
\end{enumerate}

\paragraph{Result.}
The paper maps cleanly to EBF, revealing the complementarity structure
driving housing market changes and identifying productive extensions.

\paragraph{Discussion.}
This systematic method can be applied to any new economics paper. The
five steps---finding, complementarity, context, equilibrium, restrictions---
provide a template for ongoing literature integration.


%%%%%%%%%%%%%%%%%%%%%%%%%%%%%%%%%%%%%%%%%%%%%%%%%%%%%%%%%%%%%%%%%%%%%%%%%%%%%%%
\section{Implications for Practice}
\label{sec:J-practice}
%%%%%%%%%%%%%%%%%%%%%%%%%%%%%%%%%%%%%%%%%%%%%%%%%%%%%%%%%%%%%%%%%%%%%%%%%%%%%%%

\begin{tcolorbox}[colback=green!5!white, colframe=green!75!black,
    title=Practical Implications]
\begin{enumerate}
\item \textbf{Literature monitoring:} Use EBF as lens for evaluating new
    research---identify complementarities, context, and restrictions
\item \textbf{Research design:} Frame new studies in EBF terms to connect
    to existing literature and identify gaps
\item \textbf{Teaching:} Use framework to unify disparate topics in
    economics courses
\item \textbf{Policy analysis:} Apply $(C, \Psi, C^*, K, Q)$ structure
    to evaluate policy proposals
\item \textbf{Cross-disciplinary work:} Use common language to bridge
    economics, psychology, sociology, and political science
\end{enumerate}
\end{tcolorbox}


%%%%%%%%%%%%%%%%%%%%%%%%%%%%%%%%%%%%%%%%%%%%%%%%%%%%%%%%%%%%%%%%%%%%%%%%%%%%%%%
%%%%%%%%%%%%%%%%%%%%%%%%%%%%%%%%%%%%%%%%%%%%%%%%%%%%%%%%%%%%%%%%%%%%%%%%%%%%%%%
%
%                    PART 4: BACK MATTER
%
%%%%%%%%%%%%%%%%%%%%%%%%%%%%%%%%%%%%%%%%%%%%%%%%%%%%%%%%%%%%%%%%%%%%%%%%%%%%%%%
%%%%%%%%%%%%%%%%%%%%%%%%%%%%%%%%%%%%%%%%%%%%%%%%%%%%%%%%%%%%%%%%%%%%%%%%%%%%%%%

%%%%%%%%%%%%%%%%%%%%%%%%%%%%%%%%%%%%%%%%%%%%%%%%%%%%%%%%%%%%%%%%%%%%%%%%%%%%%%%
\section{Summary}
\label{sec:J-summary}
%%%%%%%%%%%%%%%%%%%%%%%%%%%%%%%%%%%%%%%%%%%%%%%%%%%%%%%%%%%%%%%%%%%%%%%%%%%%%%%

\begin{tcolorbox}[colback=green!10!white, colframe=green!75!black,
    title=\textbf{Appendix REC: Summary}]

\textbf{Core Question:} Can EBF characterize papers published after its
development?

\textbf{Main Contribution:}
\begin{itemize}[nosep]
    \item Prospective validation: 10/10 papers (2024-2025) successfully mapped
    \item Integration method: Five-step process for new papers
    \item Restriction identification: Framework reveals modeling choices
    \item Cross-domain coverage: Social, climate, finance, trade, gender, digital
\end{itemize}

\textbf{Key Outputs:}
\begin{itemize}[nosep]
    \item Validated framework generality
    \item Identified research extension opportunities
    \item Demonstrated ongoing integration capability
\end{itemize}

\textbf{Integration:} This appendix connects to Appendix NOB (LIT-NOBEL) as
the prospective counterpart to retrospective analysis, validating framework
applicability to new research.

\end{tcolorbox}


%%%%%%%%%%%%%%%%%%%%%%%%%%%%%%%%%%%%%%%%%%%%%%%%%%%%%%%%%%%%%%%%%%%%%%%%%%%%%%%
\section{Glossary of Symbols}
\label{sec:J-glossary}
%%%%%%%%%%%%%%%%%%%%%%%%%%%%%%%%%%%%%%%%%%%%%%%%%%%%%%%%%%%%%%%%%%%%%%%%%%%%%%%

\begin{tcolorbox}[colback=blue!5!white, colframe=blue!75!black]
\textbf{Central Reference:} For complete symbol definitions, see
\textbf{Appendix GLS} (\textit{Glossary and Notation}, \S\ref{app:glossary}).

This section lists only symbols introduced or specialized in this appendix.
\end{tcolorbox}

\vspace{0.5em}

\begin{table}[htbp]
\centering
\caption{Symbols Introduced in Appendix REC}
\label{tab:J-symbols}
\renewcommand{\arraystretch}{1.3}
\begin{tabular}{@{}clp{6cm}c@{}}
\toprule
\textbf{Symbol} & \textbf{Name} & \textbf{Definition} & \textbf{Also in G?} \\
\midrule
$C^{SS}$ & Social-social complementarity & Cross-class connection effects & \checkmark \\
$C^{Eco,F}$ & Eco-financial complementarity & Environment-income interaction & NEW \\
$C^{FP}$ & Financial-personal capital & Money-skill complementarity & \checkmark \\
$\Psi_{neighborhood}$ & Neighborhood context & Local social environment & \checkmark \\
$\Psi_{inst}^{extractive}$ & Extractive institutions & Mafia/corruption context & NEW \\
$U^D$ & Digital welfare & Satisfaction from digital activity & \checkmark \\
\bottomrule
\end{tabular}
\end{table}


%%%%%%%%%%%%%%%%%%%%%%%%%%%%%%%%%%%%%%%%%%%%%%%%%%%%%%%%%%%%%%%%%%%%%%%%%%%%%%%
\section{Critical Foundations and Objections}
\label{sec:J-foundations}
%%%%%%%%%%%%%%%%%%%%%%%%%%%%%%%%%%%%%%%%%%%%%%%%%%%%%%%%%%%%%%%%%%%%%%%%%%%%%%%

\subsection{Objection 1: Selection Bias in Papers}

\textbf{Objection:} The ten papers were selected to map well to EBF.
A random sample might fail.

\textbf{Response:} Papers were selected for topical diversity (social capital,
climate, banking, trade, institutions, technology, gender, digital) and
publication venue (top journals). The selection criterion was coverage of
different domains, not fit to framework. Moreover, we challenge critics to
identify papers that \emph{cannot} be characterized in $(C, \Psi, C^*, K, Q)$
terms.


\subsection{Objection 2: Trivial Relabeling}

\textbf{Objection:} The integration is merely relabeling existing concepts
without adding insight.

\textbf{Response:} The framework does more than relabel:
\begin{enumerate}
    \item It reveals \textbf{connections} across papers (e.g., Chetty's
          social capital and Goldin's gender both involve $\Psi_{norm}$)
    \item It identifies \textbf{restrictions}---what each paper assumes
          rather than models
    \item It suggests \textbf{extensions}---productive research directions
    \item It provides \textbf{common language} for interdisciplinary work
\end{enumerate}


%%%%%%%%%%%%%%%%%%%%%%%%%%%%%%%%%%%%%%%%%%%%%%%%%%%%%%%%%%%%%%%%%%%%%%%%%%%%%%%
\section{Open Issues and Research Agenda}
\label{sec:J-open}
%%%%%%%%%%%%%%%%%%%%%%%%%%%%%%%%%%%%%%%%%%%%%%%%%%%%%%%%%%%%%%%%%%%%%%%%%%%%%%%

\subsection{Methodological Priorities}

\begin{enumerate}
    \item \textbf{Automated integration:} Develop NLP tools to extract
        $(C, \Psi, C^*, K, Q)$ elements from paper abstracts
    \item \textbf{Completeness testing:} Systematically test framework
        against random paper samples
\end{enumerate}

\subsection{Empirical Priorities}

\begin{enumerate}
    \item \textbf{Continuous monitoring:} Track 2026+ papers for integration
    \item \textbf{Field coverage:} Extend to development, health, education
        economics
    \item \textbf{Cross-disciplinary:} Test integration with psychology,
        sociology, political science papers
\end{enumerate}


%%%%%%%%%%%%%%%%%%%%%%%%%%%%%%%%%%%%%%%%%%%%%%%%%%%%%%%%%%%%%%%%%%%%%%%%%%%%%%%
\section{References}
\label{sec:J-references}
%%%%%%%%%%%%%%%%%%%%%%%%%%%%%%%%%%%%%%%%%%%%%%%%%%%%%%%%%%%%%%%%%%%%%%%%%%%%%%%

\begin{tcolorbox}[colback=blue!5!white, colframe=blue!75!black]
\textbf{Master Bibliography:} For complete reference details, see
\texttt{bcm2\_master\_references.bib} in the repository.

This section lists appendix-specific references and data sources.
\end{tcolorbox}

\vspace{0.5em}

\subsection{Key References for This Appendix}

The following works are analyzed in this appendix:

\begin{itemize}[nosep]
    \item \textbf{Social Capital:} Chetty et al. (2022, 2024)
    \item \textbf{Climate:} Dechezleprêtre et al. (2025)
    \item \textbf{Banking:} Amador \& Bianchi (2024)
    \item \textbf{Trade:} Choi et al. (2024)
    \item \textbf{Institutions:} Fenizia \& Saggio (2024)
    \item \textbf{Market Structure:} Kwon et al. (2024)
    \item \textbf{Human Capital:} Herbst \& Hendren (2024)
    \item \textbf{Gender:} Goldin (2024)
\end{itemize}

\subsection{Appendix-Specific References}

\begin{thebibliography}{99}

\bibitem{chetty2024} Chetty, R., Jackson, M.O., Kuchler, T., Stroebel, J. et al. (2024). Social Capital and Economic Mobility. \textit{Quarterly Journal of Economics}.

\bibitem{dechezlepretre2025} Dechezleprêtre, A., Fabre, A., Kruse, T., Planterose, B., Sanchez-Chico, A., \& Stantcheva, S. (2025). Fighting Climate Change: International Attitudes. \textit{American Economic Review}, 115(4).

\bibitem{amador2024} Amador, M. \& Bianchi, J. (2024). Bank Runs, Fragility, and Credit Easing. \textit{American Economic Review}, 114(7), 2073-2110.

\bibitem{goldin2024} Goldin, C. (2024). Why Women Won. \textit{American Economic Review}, 114(6), 1515-1539.

\bibitem{fenizia2024} Fenizia, A. \& Saggio, R. (2024). Organized Crime and Economic Growth. \textit{American Economic Review}, 114(7), 2171-2200.

\end{thebibliography}

\subsection{Data Sources}

\begin{itemize}[nosep]
    \item \textbf{Paper Selection:} AER, QJE, JPE 2024-2025 issues
    \item \textbf{Working Papers:} NBER, CEPR 2024-2025 releases
\end{itemize}


%%%%%%%%%%%%%%%%%%%%%%%%%%%%%%%%%%%%%%%%%%%%%%%%%%%%%%%%%%%%%%%%%%%%%%%%%%%%%%%
% CROSS-REFERENCE FOOTER (Required)
%%%%%%%%%%%%%%%%%%%%%%%%%%%%%%%%%%%%%%%%%%%%%%%%%%%%%%%%%%%%%%%%%%%%%%%%%%%%%%%

\vspace{1em}
\hrule
\vspace{0.5em}

\noindent\textit{Cross-references:}
Appendix GLS (Central Glossary),
Appendix NOB (LIT-NOBEL: Retrospective Integration),
Appendix HOW (CORE-HOW: Complementarity),
Appendix CTW (CORE-WHEN: Context),
Appendix WAT (CORE-WHAT: Welfare Dimensions).

\noindent\textit{Master Bibliography:} \texttt{bcm2\_master\_references.bib}


% =============================================================================
% END OF APPENDIX J
% =============================================================================
